% Arbre de Diane : fil de cuivre dans une solution de nitrate d'argent (exercice 2, chapitre 8)
\begin{tikzpicture}[scale=1.15]
\pic at (0,0) {tube={hauteur=3.0,niveau=2.3,liquide=pcBleu!30}};
\draw[pcCuivre,line width=1.6pt,line cap=round] (0,0.4) -- (0,3.45) -- (0.38,3.72);
\foreach \yy in {0.55,0.9,1.25,1.6,1.95}{%
  \draw[otInk!50,line width=0.7pt,line cap=round] (0,\yy) -- (-0.15,{\yy+0.11}) (-0.09,{\yy+0.065}) -- (-0.1,{\yy+0.2})
    (0,{\yy+0.1}) -- (0.15,{\yy+0.22}) (0.09,{\yy+0.165}) -- (0.1,{\yy+0.3});
  \fill[otInk!35] (-0.15,{\yy+0.11}) circle (0.03) (0.15,{\yy+0.22}) circle (0.03);}
\pcNote[west]{(0.05,3.15)}{(0.85,3.35)}{fil de cuivre}
\pcNote[west]{(0.15,1.85)}{(0.85,2.05)}{cristaux d'argent}
\pcNote[west]{(0.16,0.8)}{(0.85,0.75)}{solution qui bleuit}
\end{tikzpicture}
\par\smallskip\legende{Au contact du cuivre, les ions argent sont réduits : des cristaux d'argent se forment le long du fil.}
